%%
% BIThesis 研究生学位论文模板 The BIThesis Template for Graduate Thesis
% This file has no copyright assigned and is placed in the Public Domain.
%%

\begin{conclusion}

本文围绕大语言模型多智能体系统中的高效通信问题展开研究。针对多智能体协作过程中存在的上下文膨胀、重复消息传播、工具观察扩散和推理成本上升等问题，本文将通信冗余划分为语义冗余问题和通信边冗余问题，并分别从通信内容和通信结构两个层面提出优化方法。本文的核心目标并非简单减少智能体数量或压缩全部文本，而是在保持多智能体协作收益的前提下，提高有效信息密度，降低冗余Token消耗和工具调用带来的额外开销，从而改善多智能体系统的准确率--成本权衡。

针对语义冗余问题，本文提出了面向语义冗余的消息压缩框架。该框架从角色、节点和时序三个维度分析多智能体通信中的冗余来源，分别设计了角色感知压缩、节点感知聚合和分歧感知压缩模块。角色感知压缩根据接收方智能体职责筛选发送方推理内容，减少与下游任务弱相关的解释；节点感知聚合对多个前驱智能体传递的相似观点进行去重，并保留来源信息；分歧感知压缩从多轮历史中提取共识答案、候选分歧和简要依据，避免完整历史在后续轮次中持续累积。实验结果表明，该框架能够在多数知识理解、数学推理和代码生成任务中保持或提升任务准确率，并相对完全通信拓扑和已有图剪枝基线降低推理成本。其中，在三个主分析模型和六个数据集的平均结果上，本文方法相对Fullconnect的平均推理成本降低29.85\%，相对AgentPrune的平均推理成本降低14.20\%。

针对通信边冗余问题，本文提出了面向通信边冗余的图剪枝算法。该方法将多智能体系统中的空间边和时间边建模为可学习概率，并以任务正确性为收益项，将模型输出Token、工具调用次数和工具观察Token三类边级成本信号共同纳入优化目标，使通信边学习能够感知不同发送方、不同轮次和不同工具观察带来的成本差异。在此基础上，本文引入Taylor一阶重要性评估，根据边参数与梯度估计移除通信边对目标函数的影响，并周期性剪除低重要性边。实验结果表明，该方法更适用于工具观察冗余较多、通信边价值差异明显的任务。在MiniMax-M3上，本文方法在AIME、GAIA、GPQA和BigCodeBench四个主实验任务上的平均准确率高于三个对比基线，且相对AgentPrune+Tools降低了平均推理成本；在Qwen3-14B上，本文方法在主实验任务中也取得了相对AgentPrune+Tools的平均成本下降。

本文的创新点主要体现在以下三个方面。首先，本文从问题建模层面将多智能体通信冗余明确拆分为消息内容层面的语义冗余和通信图结构层面的通信边冗余，为分析大语言模型多智能体系统中的成本--质量权衡提供了更清晰的框架。其次，本文面向多智能体通信消息的特点设计了角色感知、节点感知和分歧感知的压缩机制，使消息压缩不再只依赖通用文本长度缩减，而是结合接收方职责、前驱节点差异和跨轮历史状态进行针对性筛选。最后，本文将工具调用成本引入通信边学习过程，并结合Taylor一阶重要性进行动态图剪枝，使图结构优化能够更细粒度地区分高价值通信边和低价值工具观察传播路径。

从应用价值看，本文研究能够为低成本、可扩展的多智能体系统部署提供方法支持。在企业问答、代码生成、复杂推理、自动化办公和工具增强智能体等场景中，多智能体系统通常需要多轮交互和多次外部工具调用，若通信内容和通信结构缺乏控制，系统成本会随智能体数量和交互轮次快速上升。本文提出的消息压缩框架和图剪枝算法均不依赖底层大语言模型参数更新，具有较好的可插拔性和迁移潜力，可用于降低API调用成本、缩短推理延迟，并减少无关信息对下游智能体判断的干扰。因此，本文工作具有一定的经济价值和工程应用前景，也为多智能体系统从实验环境走向实际部署提供了可参考的通信优化思路。

同时，本文工作仍存在进一步完善空间。第一，当前实验主要基于固定评测集上的完整评测结果，后续需要加入更多随机种子、显著性检验和更大规模任务，以进一步验证方法稳定性。第二，本文的压缩强度、成本权重和剪枝比例仍主要采用预设配置，未来可以结合任务难度、模型不确定性和工具返回质量进行自适应调整。第三，在代码生成和多项选择等部分任务中，单纯的消息压缩或通信边剪枝仍难以完全替代执行反馈、事实核验和工具观察质量评估，后续可进一步研究消息压缩、图剪枝、检索增强和可执行验证之间的联合优化。第四，随着多智能体系统逐渐应用于真实开放环境，通信优化还需要进一步考虑隐私保护、安全约束、失败恢复和长期记忆管理等问题。

综上所述，本文从通信内容和通信结构两个层面对大语言模型多智能体系统中的冗余通信问题进行了研究，分别提出面向语义冗余的消息压缩框架和面向通信边冗余的图剪枝算法。实验结果表明，两项工作能够在不同任务场景下改善多智能体系统的准确率--成本权衡，为构建更加高效、低成本和可扩展的多智能体推理系统提供了有益探索。

\end{conclusion}
